\section{Cover topology and the choice of differential}
\label{cv:section}

The energy identity and the topology of a normalized cover are compatible,
but they are not interchangeable. This distinction is visible without
any asymptotic limit.

\begin{proposition}[Euler characteristic of a normalized root cover]
\label{cv:euler}
Let $\Omega$ be a disk and suppose $f$ is holomorphic near
$\overline\Omega$, nonzero on its boundary, with distinct interior zeros
$a_1,\ldots,a_k$ of orders $m_1,\ldots,m_k$. On the normalization $X_q$
of $w^q=f(z)$,
\begin{equation}\label{cv:generalchi}
 \chi(X_q)=q-\sum_{j=1}^k\bigl(q-\gcd(q,m_j)\bigr).
\end{equation}
For $q=2$, writing $B_{\rm odd}=\#\{j:m_j\text{ is odd}\}$ and
$N=\sum_jm_j$, this becomes
\begin{equation}\label{cv:chi}
 \chi(X_2)=2-B_{\rm odd},\qquad
 \chi(X_2)-(2-N)=2\sum_j\lfloor m_j/2\rfloor.
\end{equation}
\end{proposition}
\begin{proof}
The local normalization described in Proposition~\ref{en:rootcover}
has $d_j=\gcd(q,m_j)$ points above $a_j$, each with ramification index
$q/d_j$. Triangulate the closed disk with the zeros as vertices and with
no other branch values on edges. Each open edge and face has $q$ lifts;
ordinary vertices have $q$ lifts, while $a_j$ has $d_j$. Thus the lifted
Euler characteristic is $q\chi(\Omega)-\sum_j(q-d_j)$, proving
\eqref{cv:generalchi}. For $q=2$, $d_j$ is one for odd $m_j$ and two
for even $m_j$. Subtracting $2-N$ gives
$N-B_{\rm odd}=2\sum_j\lfloor m_j/2\rfloor$.
\end{proof}

The proof includes disconnected covers, for example the two components
of the normalized curve $w^2=z^2$. In a proper analytic family with fixed
complex zero count $N$, any regular fiber with $N$ simple zeros has
$\chi(X_2)=2-N$. Before and after a collision the regular complex fibers
therefore have the same Euler characteristic. The excess in
\eqref{cv:chi} occurs at the \emph{normalized singular fiber}, not as a
permanent difference between the two regular fibers. The unnormalized
singular curve is a different topological space.

\begin{proposition}[Reciprocal differential at a zero]\label{cv:poles}
On the normalized double cover $w^2=f(z)$, the differential
$\omega=dz/w$ is regular above a simple zero. At a zero of order
$m=2r$ it has a pole of order $r$ on each of the two local components.
At a zero of order $m=2r+1$ it has a pole of order $2r$ when $r\ge1$.
Its reciprocal energy, denoted $\mathcal E(f;\Omega)$, is
\begin{equation}\label{cv:reciprocalenergy}
 \frac i2\int_{X_2}\omega\wedge\overline\omega
 =2\int_\Omega\frac{dA}{|f|},
\end{equation}
where both sides are interpreted locally or as extended integrals.
This energy is finite at a simple zero, logarithmically divergent at a
double, and power-divergent at higher multiplicity.
\end{proposition}
\begin{proof}
Write $f(z)=z^m g(z)$ after translating the zero, with $g(0)\ne0$ and
with a local holomorphic square root of $g$. For $m=2r$ the normalized
charts have $w=\pm z^r\sqrt g$, so
$\omega=\pm z^{-r}g^{-1/2}\,dz$. For $m=2r+1$, use $z=\xi^2$ and
$w=\xi^{2r+1}\sqrt{g(\xi^2)}$ to obtain
\[
 \omega=2\xi^{-2r}g(\xi^2)^{-1/2}\,d\xi.
\]
Summing the two sheetwise densities gives \eqref{cv:reciprocalenergy}.
The radial comparison $|f(z)|\asymp|z|^m$ reduces integrability to
$\int_0^\delta r^{1-m}\,dr$.
\end{proof}

By contrast, the derivative differential $\eta_2=2\,dw$ is regular at
all multiplicities. Its finite energy counts zeros only after the
localization and normalization in Theorem~\ref{en:power}. At a pure
double $f=Az^2$, the residues of $dz/w$ are
$\pm A^{-1/2}$; they are not integer residues. The integer residue of
$d\log f$ counts ordinary multiplicity, while that of a cluster
discriminant counts collision multiplicity. These three uses of
``residue'' have different normalizations and different consumers.

\begin{proposition}[Amplitude and analytic-unit dependence]\label{cv:amplitude}
Multiplication of a heat field $H$ by a nonzero constant $c$ preserves
its zeros, their trajectories and all pair distances. Nevertheless
\begin{align*}
 \omega_{cH}&=c^{-1/2}\omega_H,&
 \oint\omega_{cH}&=c^{-1/2}\oint\omega_H,\\
 \mathcal E(cH;\Omega)&=|c|^{-1}\mathcal E(H;\Omega),&
 W(cH;\Omega)&=|c|W(H;\Omega),
\end{align*}
with a consistent square-root choice and cycles identified through the natural cover isomorphism. More generally, if $H=UP$ on a
simply connected disk, with $U$ nonvanishing, the identification
$w=\sqrt U\,v$ sends the cover to $v^2=P$, but sends the physical
differential to $dz/(\sqrt U\,v)$.
\end{proposition}
\begin{proof}
The asserted transformations follow by substitution in the differentials
and energy densities. A holomorphic nonvanishing function on a simply
connected disk has a holomorphic logarithm, hence a coherent square root.
\end{proof}

Thus root positions determine the finite branched covering data, but not
the unnormalized physical periods or energies without the analytic unit.
The normalized period matrix of a separately specified compact Riemann
surface is another object. The argument above concerns the physical
differential and the bordered cover, whose relative periods and boundary
terms remain part of its energy. For the canonical theta field the
amplitude is fixed by its source; the proposition is an information
test, not permission to change that normalization.
